\section{Proofs for joint original-action information}
\label{supp:jointinformation124}

\begin{proof}[Proof of Theorem~\ref{thm:jointinterpolation124}]
First suppose $K>0$ and $q\in\mathcal Q_I$. On an adjacent interval
write $t=x-x_i\in[0,d_i]$ and
\[
 r_i(t)=q'(x_i+t)-g_i-\mu t.
\]
Strong convexity and derivative Lipschitzness imply for $u<v$
$0\le r_i(v)-r_i(u)\le K(v-u)$. Its endpoint values are
$0$ and $s_i$, so $0\le s_i\le Kd_i$ and
\[
 \max\{0,s_i-K(d_i-t)\}\le r_i(t)\le\min\{Kt,s_i\}.
\]
Integrating these two ramps gives
\[
 \frac{s_i^2}{2K}\le \int_0^{d_i}r_i(t)\,dt
        =w_i\le s_i d_i-\frac{s_i^2}{2K}.
\]
This proves all adjacent constraints. The paid bounds hold by
definition; the derivative cap and its integrated bound give the boxes.

Conversely suppose all finite constraints hold. On each adjacent
interval define the lower and upper ramps
\[
 r_i^-(t)=\max\{0,s_i-K(d_i-t)\},\qquad
 r_i^+(t)=\min\{Kt,s_i\}.
\]
Both are continuous, nondecreasing and $K$-Lipschitz, with endpoint
values $0,s_i$. Their integral range is exactly the permitted interval
for $w_i$. Choose a convex combination
$r_i=(1-\theta_i)r_i^-+\theta_i r_i^+$ whose integral is $w_i$.
If their integrals coincide use either one; otherwise
$\theta_i$ is the ratio of the requested excess integral to the
positive difference of the two integrals and belongs to $[0,1]$.
Set
\[
 q(x_i+t)=f_i+g_i t+\frac\mu2t^2+\int_0^t r_i(s)\,ds.
\]
The values at both ends equal $f_i,f_{i+1}$, and the derivatives
at both ends equal $g_i,g_{i+1}$. Hence the pieces join to a
continuously differentiable function on the entire interval. Its
derivative increments lie in $[\mu(v-u),L(v-u)]$ within each
piece and, by addition, across the joins. Integration proves
$\mu$-strong convexity, while the upper increment bound proves
derivative Lipschitzness. Its derivative lies between the two endpoint
derivatives on every interval, so the cap is preserved. It satisfies
all measured data and realizes the required list.

If $K=0$, the derivative increment must be exactly $\mu d_i$;
integration gives $s_i=w_i=0$. Conversely those equalities join
pieces of a single quadratic with curvature $\mu$. No division by
$L-\mu$ is performed. Duplicate observations are constraints on one
node and impose no zero-length interval. Endpoints are already nodes,
so no boundary extension is missing.

Finally $\mathcal P_I(a,b)$ is a nonempty closed subset of a finite
box, hence compact. The linear contrast objective has attained extrema.
Necessity and the constructive sufficiency identify these extrema with
the original function-class extrema, proving
\eqref{eq:exactcontrast124}.
\end{proof}

\begin{proof}[Proof of Theorem~\ref{thm:dualcertificate124}]
Nonnegative combinations of the convex constraints, added to the
linear objective, make $\mathcal L$ convex. Therefore for every $z$
\[
 \mathcal L(z)\ge\mathcal L(z_0)+r^{\mathsf T}(z-z_0).
\]
For any feasible $z$, $\mathcal L(z)\le H(z)$ because
$y_j\ge0$ and $c_j(z)\le0$. Minimizing the affine expression
over the enclosing box gives exactly $B$ in
\eqref{eq:rationaldual124}. Thus $H(z)\ge B$ at every feasible
point, including the original values and derivatives. Since
$H=f_a-f_b$, this proves \eqref{eq:verifiedcontrast124}.
Replacing $H$ gives the other three conclusions by the same proof.
All derivatives and arithmetic are rational polynomials in the
specified rational inputs. The proof needs neither strong duality nor
a feasible or stationary proposed $z_0$.
\end{proof}

\begin{proof}[Policy, minimax, stopping and refinement statements]
For distinct $a,b$, the only assumed information about the implemented
future error gives $\eta(b)-\eta(a)\in[-\delta,\delta]$.
Adding this to the original contrast proves
\eqref{eq:jointpolicy124}; adding $0\le\eta(b)\le\delta$ to
the absolute bounds proves its remaining inequalities. At equal
actions the same implementation is subtracted from itself. Subtracting
the same original optimum from both policies gives
\eqref{eq:regrettransport124}. Summing actual contrasts telescopes,
so finitely many accepted changes preserve their cumulative gain.

For every admissible $q$, the function
$b\mapsto q(b)-q(a)-\mu b^2/2$ is convex. Its supremum over
$\mathcal Q_I$ is convex, proving strong convexity of $U_I$.
The uniform derivative cap makes $U_I$ $M$-Lipschitz, hence continuous,
so its minimum is attained. Also $U_I(a,a)=0$, giving $\Gamma_I\ge0$.
For every action law $\nu$, Jensen's inequality gives
\[
 \sup_q\mathbb E_\nu[q(B)-q(a)]
       \ge\sup_q\{q(\mathbb E_\nu B)-q(a)\}
       =U_I(a,\mathbb E_\nu B)\ge-\Gamma_I.
\]
The deterministic minimizer attains equality. All expectations are
finite, since the common derivative cap bounds the contrast.

For the implemented cost let $r=\Pr(B\ne a)$. When $r=0$ the
contrast is identically zero. Otherwise set
$b_*=\mathbb E[B\mid B\ne a]$ and choose the Borel error
$\eta(t)=\delta\boldsymbol1_{\{t\ne a\}}$. Conditional Jensen
and an attaining function for the fixed action $b_*$ give
\[
 \sup_{q,\eta}\mathbb E[J(B)-J(a)]
 \ge r\{U_I(a,b_*)+\delta\}
 \ge r(-\Gamma_I+\delta)
 \ge\min\{0,-\Gamma_I+\delta\}.
\]
Retention attains zero; if $\Gamma_I>\delta$, the deterministic
robust minimizer differs from $a$ and attains the other expression.
This proves \eqref{eq:jointminimax124}, including laws straddling
the incumbent and an adversary chosen before the random draw.

For the stopping ceiling, let $b^*$ minimize the true $q$ and
choose a covering point $b\in\mathcal B$ with $|b-b^*|\le h$.
Then
$q(a)-q(b^*)\le q(a)-q(b)+Mh\le k_b+Mh$, proving
\eqref{eq:jointceiling124}. No new policy can improve the current
$J(a)$ by more than $q(a)-\min q+\delta$, and the incumbent's
regret bound supplies the second ceiling. A compulsory incremental
economic cost at least their minimum cannot be recovered even by
the best feasible subsequent decision. This proves the stated
single-decision stopping dominance.

Finally the minimum of $r_k(z_k-z_{0k})$ over a box coordinate is
attained at its lower endpoint if $r_k\ge0$ and its upper endpoint
otherwise. Subtraction for two nested boxes gives
\eqref{eq:refinementcredit124}. Tightening a constraint endpoint
changes only the stated constant in $\mathcal L$; its gradient is
unchanged. This proves the additional multiplier credits.
\end{proof}

\subsection{A strict rational witness with exactly the same information}
\label{supp:jointwitness124}

\begin{proposition}[A strict joint-information gain]
\label{prop:jointwitness124}
For the rational information specified below, the joint minimax original
gain is $5/64$, while integrating the derivative cones and then
intersecting the separate endpoint-value bounds guarantees $1/16$.
The strict difference is $1/64$ under identical supplied information.
\end{proposition}

Let $A=1$, $\mu=1$, $L=2$, and supply only
\[
 q(0)=0,\quad q(1)=3/16,\qquad
 q'(0)=-1/2,\quad q'(1)=1.
\]
The incumbent is $a=0$. A cap $|q'|\le2$ is redundant for this
information. Define
$h(t)=q(t)-t^2/2+t/2$. Then $h(0)=0$, $h(1)=3/16$,
$h'(0)=0$, $h'(1)=1/2$, and $h'$ is nondecreasing and
$1$-Lipschitz. The derivative-cone upper contrast from zero is
\[
 U_0(b)=\begin{cases}
 b^2-b/2,&0\le b\le1/2,\\
 (b^2-1/4)/2,&1/2\le b\le1.
 \end{cases}
\]
The upper value derived separately from the level at one and the
lower derivative envelope is
\[
 U_1(b)=\begin{cases}
 1/16-b/2+b^2/2,&0\le b\le1/2,\\
 3/16-b+b^2,&1/2\le b\le1.
 \end{cases}
\]
Each function has global minimum $-1/16$, at $b=1/4$ and
$b=1/2$ respectively. Therefore their pointwise minimum also has
minimum $-1/16$. The derivative-envelope calculation with this
separate absolute-value intersection guarantees at most $1/16$ gain.

The joint information is stronger. At $b=3/8$ write $v=h'(b)$
and $z=h(b)$. Applying the adjacent interpolation inequalities on
both sides gives
\[
 z\le\frac38v-\frac12v^2,
 \qquad
 z\le\frac1{16}-\frac18v-\frac12v^2.
\]
If $v\le1/8$, the first right side is at most $5/128$;
if $v\ge1/8$, the second is at most $5/128$. Hence
$z\le5/128$ and
\[
 q(3/8)\le\frac5{128}+\frac9{128}-\frac{24}{128}
                  =-\frac5{64}.
\]
This bound is attained. Set
\[
 h'(t)=\begin{cases}
 t,&0\le t\le1/8,\\
 1/8,&1/8\le t\le5/8,\\
 t-1/2,&5/8\le t\le1,
 \end{cases}
 \qquad h(t)=\int_0^t h'(s)\,ds.
\]
Direct rational integration gives $h(1)=3/16$ and
$h(3/8)=5/128$. Its derivative has slopes zero or one and is
continuous. Consequently $q=h+t^2/2-t/2$ is admissible, with
curvature one or two. On the middle interval
$q'(t)=t-3/8$; on the two outside intervals its sign is respectively
negative and positive. Thus its unique minimum is at $3/8$ and is
$-5/64$. This one admissible witness rules out a uniformly guaranteed
gain larger than $5/64$ at any action. The joint calculation attains
that gain at $3/8$, so the ideal joint guarantee is exactly $5/64$,
strictly exceeding $1/16$ by $1/64$ with the same data and no new
original-value query.

With arbitrary $0\le\eta\le\delta$, the ideal guaranteed
implemented gains are respectively $\max(0,1/16-\delta)$ and
$\max(0,5/64-\delta)$. This is an exact algebraic example in the
declared smooth-convex information class. It is not an investment
sample, evidence of a trained neural advantage, or a substitute for
the independent original-economy experiment. Rational independent
arithmetic can verify every endpoint, integral, curvature and gain.
